\chapter{Összegzés  }\label{ch:summary}

Összességében úgy érzem, sikerült minden fontosabb követelményt megvalósítanom. Először dolgoztam ASP.NET Core rendszerrel, így a backend fejlesztése különösen tanulságos volt számomra. Sokat fejlődött a generikus kód írási képességem is. A szigorú alapkövetelményeknek köszönhetően könnyedén tudtam bővíteni a rendszert funkciókkal.

Lefontosabb célkitőzésem azt volt, hogy a szofter nagyon könnyedén indítható legyen helyi környezetben, minimális külső függőséggel. A szakdolgozat írása közben ismerkedtem meg jobban a GitHub, GitHub Actions és a Docker használatával. Ennek köszönhetően a kód felhőben tárolva, bizonságosan verziókezelve van. A GitHub Actions segítségével pedig egy teljesen automatizált build folyamatot hoztam létre.

Szeretnék köszönetet nyílvántítani Kovács Sándor Márknak a folyamatos támogatásért, a hasznos tanácsokért és a szakmai útmutatásért. Sok időmbe telt, mire megírtam ezt a dolgozatot, és a segítsége nélkül nem lett volna lehetséges. Köszönöm a türelmét és a biztatását páromnak, volt csoportársaimnak és barátaimnak, akik végig mellettem álltak, és segítettek megtartani a motivációmat.

\section{További fejlesztési lehetőségek}

A webalkalmazás jelenlegi állapotában is teljes funkcionalitást nyújt, azonban számos lehetőség van a további fejlesztésre. Néhány lehetséges irány:
\begin{itemize}
    \item E-mail szolgáltatás integrálása a regisztráció és jelszóvisszaállítás folyamatába.
    \item Közösségi média integráció, például Google vagy Facebook bejelentkezés.
    \item Cimkézés és kategorizálás a receptekhez, hogy könnyebben lehessen keresni és szűrni őket.
    \item Receptajánló rendszer fejlesztése, amely a felhasználói preferenciák és korábbi aktivitás alapján ajánl recepteket.
    \item Mobilalkalmazás fejlesztése a webalkalmazás mellett
    \item Keresési és szűrési funkciók továbbfejlesztése, például összetettebb szűrők
    \item Unit és integrációs tesztek bevezetése, elsősorban az egyes komponensek finomabb szintű ellenőrzése érdekében.
\end{itemize}